\lecture{1.5}{Expectation and Conditional Expectation}
\label{lec:1-5}

\subsection{Expected Value}

The expected value is the population average of a random variable.

\begin{definition}[Expected value]
For a discrete random variable $W$, the expected value, expectation, or mean is
\[
\E(W)=\sum_w w\,\Prob(W=w),
\]
where the sum is over all possible values of $W$. For a continuous random variable with density $f_W$, the expected value is
\[
\E(W)=\int w f_W(w)\,dw,
\]
provided the integral is well-defined.
\end{definition}

The discrete formula says: multiply each possible value by the proportion of the population with that value, then add across values. If $W$ is age, $\E(W)$ is the average age in the target population. If $W$ is systolic blood pressure, $\E(W)$ is the population mean systolic blood pressure.

\subsection{Expectation of Transformations}

Expectations can also be taken after transforming variables.

\begin{definition}[Expectation of a transformation]
For a function $f$, the expectation of $f(W)$ is
\[
\E\{f(W)\}=\sum_w f(w)\Prob(W=w)
\]
in the discrete case, or
\[
\E\{f(W)\}=\int f(w)f_W(w)\,dw
\]
in the continuous case.
\end{definition}

Examples include:
\[
\E(W^2), \qquad \E\{\log(W)\}, \qquad \E\{\ind(W\ge 65)\}.
\]
The last example is especially important because indicator functions convert events into numerical variables.

\subsection{Binary Variables and Probabilities}

If $Y$ is binary, with possible values 0 and 1, then its expectation equals the probability that it is 1:
\[
\E(Y)
=\sum_{y=0}^1 y\Prob(Y=y)
=0\cdot \Prob(Y=0)+1\cdot \Prob(Y=1)
=\Prob(Y=1).
\]
More generally, for any event $B$,
\[
\E\{\ind(B)\}=\Prob(B).
\]

\begin{keyidea}
Means and probabilities are not separate ideas for binary variables. The mean of a 0/1 variable is the proportion of the population with value 1.
\end{keyidea}

This fact is used constantly. If $Y$ is a disease indicator, then $\E(Y)$ is disease prevalence. If $A$ is a treatment indicator, then $\E(A)$ is the proportion treated.

\subsection{Conditional Expectation as Subgroup Mean}

Conditional expectation is the mean within a subgroup.

\begin{definition}[Conditional expectation for a discrete conditioning event]
For a discrete variable $W$ and an event $A=a$ with $\Prob(A=a)>0$,
\[
\E(W\given A=a)=\sum_w w\,\Prob(W=w\given A=a).
\]
\end{definition}

For example, if $W$ is age and $A=1$ indicates monoclonal antibody receipt, then
\[
\E(W\given A=1)
\]
is the average age among monoclonal antibody recipients. It is not the average age in the whole population.

If $Y$ is binary, then
\[
\E(Y\given A=1)=\Prob(Y=1\given A=1).
\]
Thus the conditional mean of a binary outcome is a conditional probability.

\subsection{Conditional Mean as a Function}

A conditional mean can be viewed as a function of the conditioning variables. The notation
\[
\E(W\given A)
\]
does not represent a single number until $A$ is fixed. It is a random variable whose value depends on the subgroup to which the unit belongs:
\[
\E(W\given A)
=
\begin{cases}
\E(W\given A=0), & A=0,\\
\E(W\given A=1), & A=1.
\end{cases}
\]

In contrast, $\E(W\given A=a)$ is a number once $a$ is specified.

\begin{figure}[H]
\centering
\begin{tikzpicture}[node distance=1.8cm,>=Latex]
  \node[draw, rounded corners, fill=blue!5, minimum width=3.3cm] (input) {Input subgroup\\$A=a$};
  \node[draw, rounded corners, fill=yellow!10, right=of input, minimum width=3.5cm] (machine) {Conditional mean\\function};
  \node[draw, rounded corners, fill=green!6, right=of machine, minimum width=3.5cm] (output) {Output number\\$\E(W\given A=a)$};
  \draw[->, thick] (input) -- (machine);
  \draw[->, thick] (machine) -- (output);
\end{tikzpicture}
\caption{A conditional expectation can be understood as a function that maps subgroup labels to subgroup means.}
\end{figure}

\subsection{Interpreting $\E(Y\given A,W)$}

The expression $\E(Y\given A,W)$ is a conditional mean function of both treatment and covariates. If $A$ is binary and $W$ is age, then
\[
(a,w)\mapsto \E(Y\given A=a,W=w)
\]
maps each treatment-age subgroup to the mean outcome in that subgroup.

If $Y$ is binary, then
\[
\E(Y\given A=a,W=w)=\Prob(Y=1\given A=a,W=w).
\]
For instance,
\[
\E(Y\given A=1,W=65)
\]
is the outcome prevalence among treated 65-year-olds. It should not be interpreted as the outcome that would occur if a particular person were treated. It is a descriptive mean among people actually observed with $A=1$ and $W=65$.

\begin{remark}
The expression $\E(Y\given A,W)$ is often estimated by regression. A fitted regression model is therefore best understood as an estimated conditional mean function. This interpretation is more fundamental for this course than interpreting individual regression coefficients.
\end{remark}
